\section*{Part 1}

Using the instructions provided in the project description, run memcached alone and with each iBench source of interference. The main goal of this part is to understand how memcached tail latency changes when another workload uses the same CPU/cache/memory resources.

\begin{enumerate}[label=(\alph*)]
    \item Plot a line graph with achieved QPS on the x-axis and 95th percentile latency on the y-axis.

    \textbf{Answer:} We averaged the available memcached-alone run and the interference run in \texttt{part1/}. The no-interference curve stays below 2 ms p95 until around 40K achieved QPS, and then the latency knee appears. With interference, the knee appears much earlier and p95 is already above 2 ms around 15K--20K achieved QPS. The plotted data uses achieved QPS, not target QPS.
\end{enumerate}

\begin{enumerate}[label=(\alph*)]
    \setcounter{enumi}{1}
    \item How is the tail latency and saturation point affected by each type of interference?
\end{enumerate}

\textbf{Answer:}
\begin{itemize}
    \item \texttt{ibench-cpu}: It increases p95 latency and moves the knee to lower QPS. \textit{Explanation:} memcached needs CPU cycles for request handling and network processing. CPU interference steals those cycles.
    \item \texttt{ibench-l1d}: It has a visible but smaller effect. \textit{Explanation:} memcached uses hash table lookups, so data-cache misses increase service time.
    \item \texttt{ibench-l1i}: It is usually weaker than data-cache interference. \textit{Explanation:} memcached has a small request-processing code path, so instruction-cache pressure is not the main bottleneck.
    \item \texttt{ibench-l2}: It hurts more than L1 interference. \textit{Explanation:} L2 is shared by a larger part of the working set, and misses go to slower levels.
    \item \texttt{ibench-llc}: It hurts tail latency strongly. \textit{Explanation:} evicting memcached data from the last-level cache increases many request times.
    \item \texttt{ibench-membw}: It is one of the worst interferences. \textit{Explanation:} memcached is latency sensitive, and memory-bandwidth pressure makes cache misses much slower.
\end{itemize}

\begin{enumerate}[label=(\alph*)]
    \setcounter{enumi}{2}
    \item Processor affinity:
    \begin{itemize}
        \item Explain the use of the \texttt{taskset} command in the container commands for memcached and iBench.

        \textbf{Answer:} \texttt{taskset} pins a process to selected CPU cores. In this project it lets us control if memcached and iBench share the same core or only share the memory/cache hierarchy. This makes the experiment repeatable.
        \\

        \item Why do we run some iBench benchmarks on the same core as memcached and others on a different core?

        \textbf{Answer:} CPU and L1 interference must run on the same core to stress the private resources of that core. L2, LLC, and memory-bandwidth interference can run on another core, because those resources are shared outside one core. This separation helps identify what resource is causing the slowdown.
    \end{itemize}
\end{enumerate}

\begin{enumerate}[label=(\alph*)]
    \setcounter{enumi}{3}
    \item Assuming a service level objective (SLO) for memcached of up to 2~ms 95th percentile latency at 35K QPS, which iBench source of interference can safely be collocated with memcached without violating this SLO?
\end{enumerate}

\textbf{Answer:} From the available data, memcached alone has p95 about 0.866 ms at about 35K QPS. The interference run is already much worse, so it is not safe. In general, the only safe candidates are the light L1 instruction/data cases when measurements show p95 below 2 ms at 35K. CPU, LLC, and memory-bandwidth interference should not be colocated with memcached for this SLO.

\begin{enumerate}[label=(\alph*)]
    \setcounter{enumi}{4}
    \item In the lectures you have seen queueing theory.
    \begin{itemize}
        \item Is the project experiment above an open system or a closed system? Explain why.

        \textbf{Answer:} It is an open system. \texttt{mcperf} sends requests at a target QPS, and new requests can arrive even when older requests are still waiting.
        \\

        \item What is the number of clients in the system? How did you find this number?

        \textbf{Answer:} The command uses \texttt{-T 8 -C 8 -c 8}. We interpret this as 8 client threads and 8 connections per thread, so there are 64 client connections generating load.
        \\

        \item Sketch a diagram of the queueing system modeling the experiment above. Give a brief explanation of the sketch.

        \textbf{Answer:} The model is: external arrivals $\lambda$ from \texttt{mcperf} $\rightarrow$ network queue $\rightarrow$ memcached worker/service queue $\rightarrow$ reply to client. Interference increases the service time $S$ or its variance, so waiting time and tail latency increase.
        \\

        \item Provide an expression for the average response time of this system.

        \textbf{Answer:} A simple M/M/1 approximation is $R = S/(1-\rho)$, where $S$ is average service time, $\rho=\lambda S$ is utilization, and $\lambda$ is request arrival rate. When QPS increases or interference increases $S$, $\rho$ gets close to 1 and response time rises quickly.
    \end{itemize}
\end{enumerate}

\begin{enumerate}[label=(\alph*)]
    \setcounter{enumi}{5}
    \item Cost analysis.

    \begin{itemize}
        \item What is the cost per hour for running each VM?

        \textbf{Answer:} The Part 1 cluster uses e2-highcpu style VMs from \texttt{part1.yaml}. We estimated about 0.07--0.09 USD/hour for the small client/agent/server VMs, before disk and network costs. The exact number depends on the region and current Google Cloud price.
        \\

        \item How much time did your experiments take?

        \textbf{Answer:} The actual mcperf scans are short, but cluster setup, loading data, repeated runs, and cleanup took about 1--2 hours.
        \\

        \item What is the expected total cost and actual cost?

        \textbf{Answer:} The theoretical estimate is roughly a few tens of cents for Part 1 compute time if the cluster is deleted quickly. The actual cost can be higher because disks, images, and idle time during debugging also count.
        \\

        \item Does the expected cost match the actual cost?

        \textbf{Answer:} It may not match exactly. A better estimate would include exact VM runtime from creation to deletion, persistent disk price, network traffic, region, and any time when the cluster was idle but still running.
    \end{itemize}
\end{enumerate}

\newpage
